\documentclass[10pt,a4paper]{article}
\usepackage{amsmath,amssymb,amsthm,graphicx}
\usepackage{xfrac}
\newenvironment{enumeratei}
{\bgroup\def\theenumi{\roman{enumi}}\def\theenumii{\arabic{enumii}}\begin{enumerate}}
{\end{enumerate}\egroup}

\hyphenation{addi-tion-al belong belongs denote exchan-ge exchan-ged exchan-ges}

\newcommand{\psfrac}[2]{\sfrac{(#1)}{#2}}

\newcommand\mto{\mathchoice{\longmapsto}{\mapsto}{\mapsto}{\mapsto}}
\let\oldbigsqcup\bigsqcup
\renewcommand{\bigsqcup}{\mathop{\textstyle\oldbigsqcup}\displaylimits}

\let\epsilon\varepsilon

\DeclareMathOperator{\Ess}{Ess}
\DeclareMathOperator{\Leb}{Leb}

\usepackage{tikz}
\usetikzlibrary{plotmarks}

\theoremstyle{plain}
\newtheorem{theorem}{Theorem}[section]
\newtheorem{lemma}[theorem]{Lemma}
\newtheorem{proposition}[theorem]{Proposition}
\newtheorem{corollary}[theorem]{Corollary}
\newtheorem{claim}{Claim}

\theoremstyle{definition}
\newtheorem{remark}[theorem]{Remark}
\newtheorem{example}[theorem]{Example}

\newcommand{\A}{\mathcal{A}}
\newcommand{\B}{\mathcal{B}}
\newcommand{\I}{\mathcal{I}}

\newcommand{\Z}{\mathbb{Z}}
\newcommand{\N}{\mathbb{N}}
\newcommand{\R}{\mathbb{R}}
\title{Every ergodic interval exchange transformation is minimal}
\author{Przemyslaw Berk, Frank Trujillo}
\date{doi:10.5802/jep.302; CC BY4.0}
\begin{document}
\maketitle
\noindent\textit{Editorial note.} The counted statement is the author's labelled lemma environment at source0.tex lines 321-323; its proof environment opens at line 325, the line after the blank separator, so the pairing is by adjacency and the retained proof range 325-337 is the whole environment.\par\medskip
\noindent\textit{Editorial note.} The retained range stops at the end of the proof environment at line 337; the unpaired closing brace that stands at source line 338 belongs to a group opened outside the lemma and is deliberately not retained.\par\medskip
\noindent\textit{Editorial note.} This source is new to the material-authoring work and its preamble is built from the author's own macro and theorem block (source lines 8-45) rather than from a prepared adapter.\par\medskip
\noindent\textit{Editorial note.} This article carries no label-number dump and the retained counter values are left empty for the whole source, so the printed numbering of the card is the wrapper's default; the identity of the counted item is its author label lem: minimality.\par\medskip
\noindent\textit{Editorial scope.} One original result of the pinned publisher article, counted once: ergodicity implies minimality for an interval exchange transformation (author label lem: minimality).\par\medskip
\noindent\textit{Allowed standard prerequisites (editorial).}\begin{itemize}
\item Ergodicity with respect to Lebesgue measure. The transformation preserves Lebesgue measure and has no non-trivial invariant measurable set of positive measure.. It is the hypothesis, and it is used through the full measure of the set of points of continuity.
\item Minimality. A transformation is minimal when the orbit of every point is dense.. It is the conclusion, and the proof establishes the equivalent epsilon formulation.
\item First return map to an interval. The induced transformation on an interval J is defined by the first return time of the points of J to J.. It is the object whose maximal return exponent h and discontinuity set the proof controls.
\end{itemize}\medskip
\begin{lemma}\label{lem: minimality}
If $T: I \to I$ is ergodic with respect to the Lebesgue measure then it is minimal.
\end{lemma}

\begin{proof}
We will show that for every $x,y\in I$ and every $\epsilon>0$ there exists $m\in\N$ such that $|T^{-m}y-x|<2\epsilon$. Take an interval $J:=[x,x+\epsilon)$ and consider the first return map $T_J$. Let $I^J_{\beta}$ be intervals exchanged by $T_J$ and $h_{\beta}$ the corresponding first return times. Since $T$ is ergodic, the set $\tilde I:=\bigcup_{\beta\in \mathcal B}\bigcup_{k=0}^{h_{\beta}-1} T^k(I^J_{\beta})$ is of full Lebesgue measure.

Define $h:=\max\{h_{\beta}\mid\beta\in\mathcal B\}$. Consider the set
\[
C:=\{T^{j}(\partial I_{\alpha})\mid \alpha\in\A,\ j=0,\ldots,h\}.
\]
Pick $0<\delta<\epsilon$ such that $(y,y+\delta]\cap C=\emptyset$. Since $\tilde I$ is of full measure, there exists $\tilde y\in[y,y+\delta]\cap \tilde I$. By the choice of $\delta$ and by the fact that $T$ is right-continuous, the sets $\{T^{-j}[y,\tilde y]\mid j=0,\ldots,h\}$ are a family of pairwise disjoint intervals and $T^{-1}$ acts on each of them by translation. Since $\tilde y\in\tilde I$, there exists $m\le h$ such that $T^{-m}\tilde y\in J$ and $T^{-m}[y,\tilde y]=[T^{m}y,T^{-m}\tilde y]$. Thus
\[
|T^{-m}y,x|<\epsilon+\delta<2\epsilon,
\]
which finishes the proof.
\end{proof}

\end{document}
